\clearpage\begingroup
\providecommand{\ObsSec}[2]{#1\mathord{\overset{\mathrm{s}}{,}}#2}
\providecommand{\ObsArc}[2]{#1\mathord{\overset{\prime\prime}{,}}#2}
\providecommand{\ObsJupiter}{\mathord{\tikz[baseline=-.1ex,x=1em,y=1em,line width=.045em]{
  \draw (.02,.48) .. controls (.04,.70) and (.33,.69) .. (.31,.46) .. controls (.29,.29) and (.14,.19) .. (.02,.13)--(.49,.13);
  \draw (.37,.64)--(.37,-.06);}}}
\providecommand{\ObsVesta}{\mathord{\tikz[baseline=-.1ex,x=1em,y=1em,line width=.045em]{
  \draw (.04,.02) rectangle (.40,.38);
  \draw (.22,.38)--(.22,.48) (.12,.63)--(.22,.48)--(.32,.63);}}}

\SourceLeaf{284}\HistoricalHead{279}{R}{VERSCHIEDENE BEOBACHTUNGEN IN GÖTTINGEN.}
\noindent [22.]
\begin{center}\small
\rule{.78\textwidth}{.35pt}\\[.2em]
Astronomische Nachrichten, Band XV, Nr. 345, S. 167.\\[-.2em]
\rule{.78\textwidth}{.35pt}
\end{center}
\noindent\hspace*{\parindent}\{Sternbedeckung in Göttingen\textsuperscript{[*)]}. — Herr Hofrath \textsc{Gauss} hat mir folgende von ihm beobachtete Bedeckung des Sterns $57\,\delta$ Arietis mitgetheilt:
\begin{center}\begin{tabular}{llrrrl}
1837 December 9.&Eintritt&$5^{\mathrm h}$&$2^{\mathrm m}$&$\ObsSec{57}0$&M. Z.\\
&Austritt&5&20&$13{,}0$.&
\end{tabular}\end{center}
\noindent\hspace*{\parindent}Den Eintritt hält er für sehr gut, den Austritt für unzuverlässig.\}
\begin{center}\small
\rule{.78\textwidth}{.35pt}\\[.2em]
Astronomische Nachrichten, Band XXI, Nr. 494, S. 221.\\[-.2em]
\rule{.78\textwidth}{.35pt}
\end{center}
\begin{center}[23.]\quad (\textsc{Gauss} an \textsc{Schumacher}) Göttingen 1844 Januar 2.\end{center}
\noindent\hspace*{\parindent}Den Cometen [1843 III] habe ich am 1sten December zuerst gesehen und an mehreren folgenden Abenden am Kreismikrometer beobachtet; die Beobachtungen sind aber noch nicht reducirt, da ich die Durchmesser der Mikrometerringe noch nicht scharf bestimmt habe. Bloss die letzte vom 13ten December kann ich mittheilen, da die bei jenen Durchmessern zurückbleibende Ungewissheit auf die Declinationsbestimmung geringern Einfluss hat. Zur Vergleichung diente ein Stern 9. Grösse aus \textsc{Bessels} Zone 47, dessen scheinbare Stellung für diesen Abend Herr Doctor \textsc{Goldschmidt} auf
\begin{center}$78^0\,38'\,51''\,6\qquad+3^0\,40'\,12''\,2$\end{center}
\noindent berechnet hat. Der Comet folgt
\begin{center}$9^{\mathrm h}43^{\mathrm m}13^{\mathrm s}\,7$ m. Z. um $\ObsSec{25}{86}$,\end{center}
\noindent also
\begin{center}AR $=78^0\,45'\,19''\,5$\end{center}
\noindent und war
\begin{center}$9^{\mathrm h}35^{\mathrm m}50^{\mathrm s}\,9$ m. Z. südlicher $0'\,47''\,5$,\end{center}
\noindent also
\begin{center}Decl. $=+3^0\,39'\,24''\,7$\textsuperscript{[**)]}.\end{center}
\begingroup\renewcommand{\thefootnote}{[*)}\footnotetext{Vgl. den Brief an \textsc{Schumacher} vom 22. December 1837.]}\endgroup
\begingroup\renewcommand{\thefootnote}{[**)}\footnotetext{Die schärfere Reduction dieser Beobachtung siehe Astron. Nachr., Band XXI, S. 238 (Werke VI, S. 465), wobei aber der Abdruck Werke VI einen Druckfehler (1842 statt 1843) enthält.]}\endgroup
\clearpage\SourceLeaf{285}\HistoricalHead{280}{L}{PRAKTISCHE UND SPHÄRISCHE ASTRONOMIE. BRIEFWECHSEL.}
\begin{center}\small
\rule{.96\textwidth}{.35pt}\\[.2em]
Briefwechsel zwischen \textsc{Gauss} und \textsc{Schumacher}, Altona 1860—65, IV. S. 208, V. S. 228, 247.\\[-.2em]
\rule{.96\textwidth}{.35pt}
\end{center}
\begin{center}[24.]\quad \textsc{Gauss} an \textsc{Schumacher}, Göttingen, 10. Januar 1844.\end{center}
\noindent\hspace*{\parindent}Erst gestern Abend \mbox{. . . . .} habe ich den Cometen [1843 III] wiedergesehen, nachdem es fast volle vier Wochen ununterbrochen trübe gewesen war. Aber auch gestern Abend war die betreffende Stelle nur zuweilen zwischen Wolken sichtbar. Ich konnte nur eine einzige Vergleichung mit einem Stern 9. Grösse machen, der dem Cometen etwa 8 Minuten südlich folgte. Für Declination war der Durchgang unbrauchbar: die Rectascension ergab sich
\begin{center}1844 Jan. 9.\quad $8^{\mathrm h}57^{\mathrm m}19^{\mathrm s}$\quad AR. $=77^0\,5'\,\ObsArc{30}8$\end{center}
\noindent etwa $1\frac12$ Minuten mehr als die Elemente. Nach einem vollen Monat Zwischenzeit nach den letzten den Elementen zum Grunde liegenden Beobachtungen ist dies alles, was man erwarten kann. Auch bei der Declination werden die Elemente höchstens ein paar Minuten im Fehler sein. Der Comet war \textls[120]{sehr} lichtschwach, vielleicht zum Theil Folge des ungünstigen Luftzustandes.\linebreak
\mbox{. . . . . .}
\begin{center}[25.]\quad \textsc{Gauss} an \textsc{Schumacher}, Göttingen, 4. November 1846.\end{center}
\noindent\hspace*{\parindent}Ich verfehle nicht, Ihnen \mbox{. . . . .} abermals eine Meridianbeobachtung des Transuranischen Planeten mitzutheilen:
\begin{center}\begin{tabular}{llrrr}
1846 Nov. 3.\quad $7^{\mathrm h}0^{\mathrm m}\ObsSec{50}7$ M. Z.&G. A. $=$&$327^0$&$50'$&$\ObsArc3{5}$\\
&Abweichung $=$&$-13$&34&$13{,}1$.
\end{tabular}\end{center}
\noindent\hspace*{\parindent}Die Beobachtungen stimmen noch immer gut mit einer Kreisbahn (Halbmesser $=29{,}932$); indessen lässt sich daraus noch nichts über den Grad der Kleinheit der Ellipticität schliessen.
\begin{center}[26.]\quad \textsc{Gauss} an \textsc{Schumacher}, Göttingen, 1. December 1846.\end{center}
\noindent\hspace*{\parindent}Ich schicke Ihnen \mbox{. . . . .} hier noch eine Meridianbeobachtung des neuen Planeten [Neptun], die jedoch nicht ganz so gut sein wird wie die frühern, da der Planet in der etwas dunstigen Luft nur schwer zu sehen war.
\begin{center}\begin{tabular}{llrrr}
1846 Nov. 30.\quad $5^{\mathrm h}15^{\mathrm m}\ObsSec{20}6$&G. A. $=$&$327^0$&$59'$&$\ObsArc8{7}$\\
&Abweichung $=$&$-13$&30&$52{,}4$.
\end{tabular}\end{center}
\clearpage\SourceLeaf{286}\HistoricalHead{281}{R}{VERSCHIEDENE BEOBACHTUNGEN IN GÖTTINGEN.}
\begin{center}\textbf{BEMERKUNGEN.}\end{center}
\begingroup\small
\noindent\hspace*{\parindent}In dem im Nachlass befindlichen \textit{Tagebuch der Sternwarte}, das in zwei Bänden die Zeit vom 3. November 1808 bis zum 8. Oktober 1818 und vom 30. Oktober 1820 bis zum 15. April 1822, jedoch mit Unterbrechungen, umfasst, sind anfangs ziemlich regelmässig, später mehr vereinzelt, die auf der Sternwarte zu Göttingen angestellten Beobachtungen, häufig nebst ihrer Reduktion, eingetragen. Die Beobachtungen der ersten Zeit sind vorwiegend von \textsc{Harding}, im besonderen am Mauerquadranten, in der späteren vorwiegend von \textsc{Gauss} angestellt, teilweise auch von \textsc{Gauss}’ Schülern und Gästen. Reduziert und veröffentlicht ist davon nur ein Teil. In Nr. 1—12 sind aus dem \textit{Tagebuch} diejenigen noch unveröffentlichten Beobachtungen von \textsc{Gauss} wiedergegeben, deren Reduktion aus dem Buche entnommen werden konnte.

Beim Abdruck ist die ursprüngliche Form der Eintragung wiedergegeben worden, um einen Einblick in die Eigenart des \textit{Tagebuchs} zu gewähren. Unberücksichtigt blieben hierbei: eine Beobachtung der Ceres von 1809 August 10., eine Verfinsterung des 1. Jupitermondes vom 5. Februar 1810, eine Beobachtung des Kometen 1811 I., eine Bedeckung eines Sternes im Taurus vom 16. Dezember 1812, sowie mehrere andere, besonders am Mauerquadranten, wegen Fehlens oder Unvollständigkeit der Reduktion. Man vergleiche den Aufsatz über \textsc{Gauss}’ astronomische Arbeiten, Werke XI\textsubscript{2}.

Die Notizen 13.—26. betreffen Beobachtungen, die in den Werken bisher nicht abgedruckt sind, insbesondere solche aus dem Briefwechsel, die \textsc{Gauss} selbst nicht veröffentlicht hat.
\begin{flushright}\textsc{Brendel}.\end{flushright}
\endgroup
\begin{center}\rule{.24\textwidth}{.35pt}\end{center}
\clearpage\endgroup
